% ===== tr-01 块4：第 6 章 pp.1--4（原文页 1--4）=====
% [原文 p.1]
\section*{摘要}
本文综述作者与 A. Jacob 合作、关于有理椭圆曲面与 del Pezzo 曲面上 Strominger--Yau--Zaslow 镜像对称的近期工作。我们讨论若干应用，包括 ALH* 引力瞬子的 Torelli 定理，并说明如何用我们的结果证明所有 ALH* 引力瞬子都能紧化到一个弱 del Pezzo 曲面，从而重新得到 Hein--Sun--Viaclovsky--Zhang 的近期结果。

\section{引言}
作者与 A. Jacob 最近对由 del Pezzo 曲面与有理椭圆曲面产生的某些对数 Calabi--Yau 曲面上的 Strominger--Yau--Zaslow 镜像对称猜想的若干方面做了详细研究 [10, 11, 12]。本综述的目的是复述并总结其中一些结果，并讨论一些应用与进一步的问题。

我们首先回顾 SYZ 镜像对称猜想的一个一般表述：

\begin{conjecture}{6.1.1}[Strominger--Yau--Zaslow]
设 $(\check X,\check\omega)$ 是 Calabi--Yau 流形，$\check{\mathcal{M}}_{\mathrm{cplx}}$ 表示 $\check X$ 上复结构的模空间。则对充分接近大复结构极限的复结构 $\check J\in\check{\mathcal{M}}_{\mathrm{cplx}}$，下列结论成立：
\begin{enumerate}
\item[(1)] $(\check X,\check J,\check\omega)$ 容许一个特殊拉格朗日环面纤维化 $\check\pi:\check X\to\check B$，其底 $\check B$ 带有整仿射结构。
\item[(2)] 存在另一个 Calabi--Yau 流形 $(X,J,\omega)$，带有特殊拉格朗日纤维化 $\pi:X\to B$，且 $B$ 带有整仿射结构。
\item[(3)] 设 $\mathcal{M}_{\mathrm{Käh}}$ 表示 $X$ 的复化 Kähler 模空间。则存在一个镜像映射 $q:\check{\mathcal{M}}_{\mathrm{cplx}}\to\mathcal{M}_{\mathrm{Käh}}$，它是局部微分同胚，且满足 $\operatorname{Im}(q(\check J))=\omega$。
\item[(4)] 存在同构 $\varphi:\check B\to B$，它交换复仿射结构与辛仿射结构，并且使得位于 $\check b\in\check B$ 与 $\varphi(\check b)\in B$ 之上的特殊拉格朗日环面纤维的 Riemann 体积互为倒数。
\end{enumerate}
\end{conjecture}
% [Tutor Note: $\check{\mathcal{M}}_{\mathrm{cplx}}$ 是 $\check X$ 上复结构的模空间（花体 M 加 check，下标 cplx），与 Kähler 侧的 $\mathcal{M}_{\mathrm{Käh}}$ 相对；不要把花体 M 抄成普通 $M$。]

% [原文 p.2]
SYZ 猜想的这个表述或许是该猜想最乐观、最具雄心的版本。注意：若 $X$ 是 $n$ 维紧 Calabi--Yau 流形，则镜像 $\check X$ 也是 $n$ 维紧 Calabi--Yau 流形，而第三点蕴含 Hodge 数的互换
\[
h^{n-1,1}(\check X)=h^{1,1}(X).
\]
% [Tutor Note: 这里的 Hodge 数互换是第 (3) 点（镜像映射 + 复化 Kähler 模空间）在紧 Calabi--Yau 情形下的推论；注意指标是 $h^{n-1,1}$（不是 $h^{n-2,1}$），对曲面即 $h^{1,1}$ 互换。]

对我们而言一个重要而带有哲学意味的要点是：空间 $\mathcal{M}_{\mathrm{Käh}}$ 应当被看作 $X$ 上 Calabi--Yau 结构的（复化）辛模空间，而 $\check{\mathcal{M}}_{\mathrm{cplx}}$ 应当被看作 $\check X$ 上 Calabi--Yau 结构的复模空间。这种用 Calabi--Yau 结构重新表述模空间的做法，归功于 Yau 对 Calabi 猜想的解决 [64] 以及 Bogomolov--Tian--Todorov 定理 [4, 61, 63]。最后我们指出，Hitchin [33] 的工作表明：猜想 6.1.1 中特殊拉格朗日环面纤维化底上的整仿射结构继承自 $X,\check X$ 的复/辛几何；并且猜想 6.1.1 的第 (4) 点可以看作 SYZ 猜想的核心原则------镜像对称就是 T 对偶------的一个推论。

作者与 A. Jacob 最近对某些完备非紧 Calabi--Yau 流形证明了猜想 6.1.1 [10, 11]。我们非正式地概述这些结果。首先介绍与本文相关的对象：
\begin{itemize}
\item 设 $Y$ 是 del Pezzo 曲面，$K_Y^2=d$，其中 $0<d\le 9$，并设 $D\in|-K_Y|$ 是一条光滑椭圆曲线。那么 Tian--Yau 的一个奠基性结果 [62] 表明：$X=Y\setminus D$ 容许一个恰当的、完备的 Kähler Ricci 平坦度量 $\omega_{TY}=\sqrt{-1}\partial\bar\partial\varphi_{TY}$，它渐近于 \textbf{Calabi ansatz}。
\item 设 $\check Y\to\PP^1$ 是有理椭圆曲面（见定义 6.2.2），并设 $\check D\in|-K_Y|$ 是 Kodaira $I_d$ 型奇异纤维。那么 Hein 的一个结果 [32] 表明：$\check X=\check Y\setminus\check D$ 容许许多完备的 Kähler Ricci 平坦度量，它们渐近于 \textbf{半平坦 Calabi--Yau 度量}。
\end{itemize}
% [Tutor Note: 第二条中原文印作 $\check D\in|-K_Y|$，但按上下文 $\check Y$ 才是该有理椭圆曲面，这里应为 $|-K_{\check Y}|$；900 dpi 裁图 work/tmp/z6p02c.png 显示 $K$ 的下标为无 check 的 $Y$，是原文笔误。正文照原样保留。]
% [Tutor Note: $\textbf{Calabi ansatz}$ 与 $\textbf{半平坦 Calabi--Yau 度量}$ 在原文中为粗体（术语强调），译文相应保留粗体。]

为了证明猜想 6.1.1，需要处理若干个几何问题：
\begin{enumerate}
\item[(1)] 非紧 Calabi--Yau 空间 $X,\check X$ 是否容许特殊拉格朗日纤维化？
\item[(2)] 完备 Calabi--Yau 度量没有一般的唯一性定理，那么 Kähler 模空间 $\mathcal{M}_{\mathrm{Käh}}$ 指的是什么？能否构造一个有限维模空间来参数化完备、Kähler、Ricci 平坦的度量？
\item[(3)] 假设我们能定义 $\mathcal{M}_{\mathrm{Käh}}$，如何定义从 $X$ 的复模空间到 $\check X$ 的 Kähler 模空间的镜像映射？
\end{enumerate}

% [原文 p.3]
粗略地说，[10] 处理问题 (1)，而 [11] 处理问题 (2) 与 (3)。在本综述中，我们将复述这些结果，解释一些应用，并强调一些公开问题。

\section{del Pezzo 曲面中的特殊拉格朗日纤维化}
本节我们按照作者与 A. Jacob 在 [10] 中的论证，描述 del Pezzo 曲面中特殊拉格朗日纤维化的构造。

\subsection{Calabi ansatz}
设 $D$ 是 $(n-1)$ 维射影 Calabi--Yau 流形，$K_D\simeq\mathcal{O}_D$，并设 $p:E\to D$ 是正线丛。记 $Y_{\C}$ 为 $E$ 全空间中零截面的一个管状邻域，$X_{\C}:=Y_{\C}\setminus\{0\}$ 为 $Y_{\C}$ 中零截面的补集。设 $\Omega_D$ 表示 $D$ 上的全纯体积形式。则 $X_{\C}$ 容许由下式给出的全纯体积形式
\[
\Omega_{\C}=p^*\Omega_D\wedge\frac{dw}{w},
\]
其中 $w$ 是沿 $E$ 的纤维的坐标。设 $\omega_D$ 是 $D$ 上唯一的 Ricci 平坦度量，满足 $[\omega_D]=2\pi c_1(E)$，并把 $\Omega_D$ 归一化使得
\[
\omega_D^{n-1}=\frac{1}{2}(\sqrt{-1})^{(n-1)^2}\Omega_D\wedge\overline{\Omega_D},
\]
再设 $h$ 是 $E$ 上的 Hermitian 度量，使得 $\omega_D=-\sqrt{-1}\partial\bar\partial\log h$。考虑 Kähler 形式
\[
\omega_{\C}=\sqrt{-1}\partial\bar\partial\frac{n}{n+1}\bigl(-\log|\xi|_h^2\bigr)^{\frac{n+1}{n}},
\]
其中 $\xi\in E$。由直接计算可得
\[
\omega_{\C}^n=\frac{1}{2}(\sqrt{-1})^{n^2}\Omega_{\C}\wedge\overline{\Omega_{\C}}.
\]
几何上，若用 $r(x)$ 表示 $x$ 到固定点 $x_0\in X_{\C}$ 的距离，则直接计算表明
\begin{equation}
|\nabla^k Rm|\le C_k r^{-\frac{k+2}{n+1}},\qquad \operatorname{inj}(x)\sim r^{\frac{1-n}{n+1}},\qquad \operatorname{Vol}(B_r(x_0))\sim r^{\frac{2n}{n+1}}.
\tag{6.2.1}
\end{equation}
换言之，Riemann 曲率以多项式速率衰减，但单射半径在无穷远处退化。因此 $X_{\C}$ 不具有"有界几何"。此外容易验证，该度量在 $|\xi|\to 0$ 时完备，而在 $|\xi|\to 1$ 时不完备。
% [Tutor Note: 记号 $X_{\C},Y_{\C},\Omega_{\C},\omega_{\C}$ 中的下标是黑体 $\C=\mathbb{C}$（表示复化/复结构下的对象），不是普通字母 C；$X_{\C}=Y_{\C}\setminus\{0\}$ 表示挖去零截面。]
% [Tutor Note: (2.1)（本译本式 (6.2.1)）中三个量都是 $r\to\infty$ 时的渐近：$|\nabla^k Rm|\le C_k r^{-(k+2)/(n+1)}$、$\operatorname{inj}(x)\sim r^{(1-n)/(n+1)}$、$\operatorname{Vol}(B_r(x_0))\sim r^{2n/(n+1)}$。指数分母都是 $n+1$，不要抄成 $n$。]

在其奠基性工作 [62] 中，Tian--Yau 利用 Calabi ansatz $(X_{\C},\omega_{\C},\Omega_{\C})$ 构造了完备非紧的 Kähler Ricci 平坦流形。确切地说，给定 $n$ 维 Fano 流形 $Y$ 与光滑反典范除子 $D$，补集 $X=Y\setminus D$ 容许一个在 $D$ 上有单极点的全纯体积形式 $\Omega$，且 $\Omega$ 在 $\C^*$-伸缩下唯一。显然 $D\subset Y$ 的一个邻域可以光滑地等同于 $D$ 的法丛中零截面的一个邻域。Tian--Yau 发展出一套稳健的工具来求解非紧复流形上的复 Monge--Ampère 方程，并应用他们的技巧构造了 $X=Y\setminus D$ 上渐近于 Calabi ansatz 的完备 Ricci 平坦度量 $\omega_{TY}$。Hein [32] 后来改进并强化了 Tian--Yau 方法，得到如下结果：在 Tian--Yau 空间中，二元组 $(\omega_{TY},\Omega)$ 收敛到 $((\Phi^{-1})^*\omega_{\C},(\Phi^{-1})^*\Omega_{\C})$，且具有指数衰减，其中 $\Phi:X_{\C}\to X$ 是到其像上的微分同胚 [36]。因此，我们期望能把 $X_{\C}$ 的几何性质转移到 $X$ 的无穷远邻域上。
% [Tutor Note: $((\Phi^{-1})^*\omega_{\C},(\Phi^{-1})^*\Omega_{\C})$ 是把 Calabi ansatz 的度量与体积形式沿 $\Phi^{-1}$ 拉回到 $X$ 的无穷远邻域；$\Phi:X_{\C}\to X$ 是到其像上的微分同胚（原文 "onto is image" 为 "onto its image" 的笔误）。]

% [原文 p.4]
\subsection{Tian--Yau 空间中的特殊拉格朗日}
我们现在解释如何在 Tian--Yau 空间中产生特殊拉格朗日。特殊拉格朗日的概念由 Harvey--Lawson [35] 在其关于校准几何的奠基性工作中引入。设 $X$ 是带有全纯体积形式 $\Omega$ 与 Ricci 平坦度量 $\omega$ 的 Calabi--Yau 流形。半维子流形 $L$ 若满足 $\omega|_L=0$ 则称为拉格朗日的。由局部坐标计算可得 $\Omega|_L=e^{i\theta}\mathrm{vol}_L$，其中 $\theta:L\to S^1$ 称为 $L$ 的相函数。若 $\theta:L\to S^1$ 可以提升为映射 $\theta:L\to\R$，则称该拉格朗日子流形是分次的。若 $\theta$ 为常数，则称 $L$ 是特殊拉格朗日的。作为 校准子流形，特殊拉格朗日在其余调类中是体积最小的，并且它们在几何与物理中起核心作用。特殊拉格朗日的典型例子来自
\begin{itemize}
\item 反全纯反辛对合的实轨迹，
\item 超 Kähler 流形中全纯拉格朗日子流形的超 Kähler 旋转 [35]，
\item 已知特殊拉格朗日的形变 [52]，
\item 辛约化中特殊拉格朗日的提升 [23]。
\end{itemize}
% [Tutor Note: 相函数 $\theta$ 满足 $\Omega|_L=e^{i\theta}\mathrm{vol}_L$；"分次（graded）"指 $\theta$ 可提升为实值函数 $\theta:L\to\R$，"特殊（特殊拉格朗日）"指 $\theta$ 为常数------两个概念不可混淆。]
% [Tutor Note: 原文 "they play a central role geometry and physics" 缺介词 in（应为 "in geometry and physics"），译文按语义补足，不改变数学内容。]

现在回到 Calabi ansatz 的设定。假设给定 $D$ 中关于 $\omega_D,\Omega_D$ 的任意特殊拉格朗日 $L\subseteq D$，相为 $\theta$，即 $\Omega_D|_L=e^{i\theta}\mathrm{vol}_L$，我们可以联系一个子流形
\[
L_\varepsilon=p^{-1}(L)\cap\{|\xi|_h=\varepsilon\}\subseteq X,
\]
由直接计算，它关于 $\omega_{\C},\Omega_{\C}$、相为 $\theta+\frac{\pi}{2}$ 也是特殊拉格朗日的。特别地，当 $D$ 是椭圆曲线或带特殊拉格朗日纤维化的 K3 曲面时\footnote{例如，任何代数 K3 曲面都容许特殊拉格朗日纤维化 [10]。}，$X_{\C}$ 也同样容许特殊拉格朗日纤维化。
% [Tutor Note: 原文此处把 $L_\varepsilon$ 的所在空间印作 $\subseteq X$，但 $p^{-1}(L)$ 位于 $E$ 的管状邻域 $Y_{\C}$ 中、且本节明确处在 Calabi ansatz 设定，故应为 $\subseteq X_{\C}$；正文照原样排。]
% [Tutor Note: 相从 $\theta$ 变为 $\theta+\frac{\pi}{2}$：这一 $\frac{\pi}{2}$ 的相位移动是超 Kähler 旋转在纤维方向上的效果，不是笔误。]

利用 Tian--Yau 空间渐近于 Calabi ansatz 这一事实 [36]，人们可以对 $\Phi(L_\varepsilon)$ 施行 Moser 技巧，从而在 $X$ 的无穷远处附近产生一个拉格朗日子流形 $M_\varepsilon$，其相函数为 $O(e^{-C/\varepsilon})$ 阶，其中 $C$ 是某个正常数。特别地，$M_\varepsilon$ 是分次拉格朗日的。想法是把 $M_\varepsilon$ 扰动成 Tian--Yau 空间中真正的特殊拉格朗日。处理这一扰动问题的一种途径是使用拉格朗日平均曲率流来形变 $M_\varepsilon$，我们将在下面解释。

一般而言，若 $M$ 是 Calabi--Yau 流形中的分次拉格朗日，则平均曲率由 $\vec H=J\nabla\theta$ 给出，其中 $J$ 是殆复结构，$\nabla\theta$ 是相函数的 Riemann 梯度。
% [Tutor Note: 分次拉格朗日的平均曲率向量为 $\vec H=J\nabla\theta$（$\vec H$ 带箭头，表示法向平均曲率向量）；当 $\theta$ 为常数时 $\vec H=0$，即该拉格朗日是极小的，这正是"特殊拉格朗日"与"平均曲率流不动点"之间的联系。]
